% Transcription — fonds Alexandre Grothendieck, shelfmark 156-6, batch 5 (pages 81-93).
% Established from the facsimile archives/batches/156-6/batch-05.pdf (a copy of
% 156-6.pdf fetched through the project's relay; no cover sheet in the batch,
% so PDF page k = folder page 80+k), per the skill transcribe-grothendieck. The
% folder is ninety-three pages in five batches; this is the fifth and last,
% thirteen pages. Batches 1, 2 and 4 of this folder and batches 1, 2 and 5 of
% the first draft (folder 156-4, GF IV) were read for notation; batch 3 (pages
% 41-60) did not exist when this batch was written. Page 81 opens a new
% subject after the « saturation » remark closing batch 4 (page 80).
%
% Pass: Opus 5.5 (claude-opus-5-5), 2026-09-26 — first pass, unchecked against the pages by a human.
%
% Pages skipped: none. Pages summarised: none. All thirteen leaves are his
% mathematics, fast drafting in ink, and all are transcribed; page 93 holds
% a dozen lines and ends on a short rule, the end of the folder.
%
% His pagination 81-93 stands at the head of each leaf and coincides with the
% archivists' (noted on page 81). One date in his hand: « 22. Juin »,
% underlined, in the left margin of page 87 against « J'ai l'impression… » —
% two days after the 18-20/06/1986 of the catalogue title, a day after the
% « 21. juin » of batch 4; transcribed as a \marginal with a note.
%
% Numbered runs, recorded where they stand: the atelier axioms At 13 (p. 83),
% At 14 (p. 86), At (fort) 14 bis (p. 87), At (fort) 14 ter (p. 87), all
% continuing batch 4's At 9-At 12; he cites At 4, At 9, At 9', At 9_Σ, At 10,
% At 11, At 11 bis, At 12 (pp. 81-93). For ensemblist ateliers, « Atens »
% axioms: Atens 2, Atens* 1-3, Atens* 4 (p. 89, his star read as in batch 4),
% Atens 2_Σ = Atens* 4_Σ (p. 89), Atens 4 (p. 90), Atens fort 5 (p. 91); the
% label of the Sup-filtrant axiom on page 90 is written over a struck number.
% Sub-ateliers: « Sous at 1, 2 » cited « p. 47 » (p. 91, batch 3, the second
% digit doubtful), Sous-at (épais) 3, Sous-at (quasi-dense) 3', Sous-at 4
% (p. 92). A three-item programme 1°)-3°) on page 88. He cites (1.58) « page
% 59 » (p. 81), « le lemme p. 82 » (p. 85) and « (p. 55-57) » (p. 88, the
% first number doubtful), all as plain page numbers.
%
% What the batch is about. Pages 81-86: subdivisions. A Proposition (p. 81):
% if F' ≪ F and some non-empty G is disjoint from F' but not from F, then one
% such G can be taken ≪ F; reduced via At 9 to lieux x, y and a Lemme
% (p. 82: for F' ≪ F, G ≤ F and the induced F'_G ≪ G, H ≪ G is disjoint from
% F' as soon as it is disjoint from F'_G); the axiom At 13 (p. 83) and a
% Corollaire listing four equivalent conditions on F' ≪ F ((i) every G
% disjoint from F' is disjoint from F, (ii) no lieu of omb F is disjoint from
% F', (iii) F' maximal among refinements of F for ≤, (iv) for ≪, cancelled),
% whose proof on page 84 ends on « Mais impliquent-elles (iv) ? ». Page 85:
% Definition — F' is a « subdivision » of F, F' ≼ F, when these hold; ≼ is
% an order; subdivisions restrict to sub-figures, and are detected on a
% covering family F = Sup F_i (Corollaire, proof p. 86). Pages 86-88:
% « recollement des subdivisions », At 14 (a subdivision G' of G ≤ F extends
% to F), At (fort) 14 bis (canonical extension, with a cartesian square of
% ordered sets), and, dated 22 June, At (fort) 14 ter for X ∈ 𝓜 and F' ≼ ∂X
% (unique X' ≼ X with ∂X' = F'), false for simplicial ateliers (a small
% triangle is drawn as counter-example). Page 88: a programme (make the
% axioms explicit for ensemblist ateliers; their stability under
% sub-ateliers; the « saturation » of an atelier). Pages 88-91: the axioms
% for an ensemblist atelier (ℒ, 𝔉 ⊂ Fig(ℒ)), which satisfies At 1-At 8
% automatically: At 9 becomes Atens 2 (X, Y with disjoint supports are
% compatible), At 9_Σ becomes Atens 2_Σ, At 10 is automatic, the Sup-filtrant
% axiom is empty for figures of finite type, At 12 becomes Atens 4, At 13
% always holds, At 14 bis becomes Atens fort 5 — « pas vrai pour les ateliers
% simpliciaux ! ». Pages 91-93: sub-ateliers 𝔉' ⊂ 𝔉 — At 1-At 7 pass to 𝔉';
% ℒ ∩ 𝔉' ⊂ ℒ'; « épais » and « quasi-dense » sub-ateliers (Sous-at 3, 3'),
% Sous-at 4 (disjointness in 𝔉 persists in 𝔉'), and which of At 9-At 14
% then pass to 𝔉'. The folder ends there.
%
% Relation to the first draft. GF IV takes ≼ (subdivision) as a primitive
% order with ≪ derived (156-4 batch 5, p. 84) and states « recollement des
% subdivisions » as its axiom C8 (GF IV p. 86); here ≼ is defined from ≪ and
% disjointness (p. 85) and the gluing is Corollaire 2 / At 14 (p. 86). Noted
% at pages 85 and 86.
%
% Reading hazards fixed once here, aligned with batches 2 and 4: 𝔉
% \mathfrak{F}, 𝓜 \mathcal{M}, ℒ \mathcal{L}; disjointness ∥ \parallel;
% compatibility (crossed chevrons) \between, with 𝓜 or 𝔉' set under it as he
% writes it; interior refinement \overset{\circ}{\ll}; « omb », « Omb »,
% « Fig » in roman. New in this batch: ∥ with a bar over it, « non
% disjoint », set \overline{\parallel}; his ≤ with a wavy lower stroke,
% subdivision, set \preccurlyeq as in batch 1 (and 156-4); ∂ for the boundary
% of a multistrate; « OPS » (on peut supposer) kept as written. The small
% vertical schemes of pages 81-83, 85 and 86 (≤, ≪ written sideways between
% stacked terms) are given as lists of relations with a note; the square of
% page 87 is set in tikz-cd; the little triangle of page 88 is described.
% His underlining is set \emph.
%
% Hard places, read only in part: the struck passage and its box on page 82;
% the long slanted left margins of pages 83, 85 and 88; the tangled
% parenthesis on simplicial ateliers (p. 88); several connecting words on
% pages 90-93 (one word, « discordante » (?), recurs on pages 91 and 93).
% Variances kept and flagged: « omb(X) » in (ii) of the Corollaire of
% page 83 where the proof (p. 84) uses omb(F); a slashed ≤ for ≠ (p. 84);
% « x ∥ F' » without the bar at the end of page 86's proof; « X ≬ Y » for
% X' ≬ Y' at the end of Atens 4 (p. 90); the direction 𝔉' → 𝔉 of the At 10
% sentence on page 93.
\documentclass[11pt,a4paper]{article}
\input{../preamble/grothendieck}

\folder{156-6}
\batch{5}
\pages{81}{93}
\dating{1986}
\watermark{Édition de démonstration}
\foldertitle{[Chapitre] VI. Analysis situs (deuxième mouture) : notes manuscrites (18-20/06/1986)}

\begin{document}

\page{81}
\note{pagination de sa main 81 à 93 en tête des feuillets, qui coïncide avec celle des archivistes ; sous le numéro, un court trait horizontal sépare ce qui suit de la fin de la page 80}

J'en viens aux \struck{axiomes} \emph{subdivisions}, dont il n'avait pas encore été question jusqu'à présent ! J'ai envie de poser

\textbf{\struck{\ill{}} Proposition}\note{« Proposition » est écrit au-dessus d'un mot souligné puis lourdement biffé (peut-être « Axiome »)} Soient $F$ une figure, $F' \ll F$. Supposons qu'il existe une figure $G$ non vide, avec $G \parallel F'$, $G \mathrel{\overline{\parallel}} F$. Alors il existe \struck{une} une telle figure avec $G \ll F$.
\note{le signe $\overline{\parallel}$, deux traits verticaux surmontés d'une barre, s'oppose ici à $\parallel$ (disjonction) : « non disjoint ». Il revient pages 82 à 86}

\struck{En vertu de} On prendra dans \uncertain{premier} pour $G$ une $X \in \mathcal{M}$, et \struck{une} une $x \in \mathcal{L}$, \struck{pour} en prenant n'importe quel $X$ dans \uncertain{$\mathrm{Omb}(G)$} avec \struck{\ill{}} $x$ dans $\mathrm{Omb}(G)$. (NB \ill{} la vérification de (1.58 page 59))\note{la formule (1.58) et sa page 59 sont au lot 3 ; le lot 4 les cite de même}

Pour trouver une énoncé minimal qui implique la prop., notons que dans la situation initiale, $G$, l'hypothèse $G \mathrel{\overline{\parallel}} F$ signifie aussi (grâce à At 9) qu'il existe \add{$x \in \mathrm{Omb}(F)$,} $y \in \mathrm{Omb}\, G$, tels que \struck{\ill{}} $x \mathrel{\overline{\parallel}} y$, et même, prenant \ill{}, $y \parallel F'$ ! Soit donc $X \in \widetilde{F}$ tel que $x \in \mathrm{Omb}(X)^{\circ}$ i.e. $x \mathrel{\overset{\circ}{\ll}} X$.\note{« une énoncé » est sur la page ; « $\in \widetilde{F}$ » : le tilde est lu sous le $F$ en fin de ligne, douteux}

\[
  F' \ll F, \qquad F'_X \leq F', \qquad X \leq F, \qquad F'_X \ll X, \qquad x \mathrel{\overset{\circ}{\ll}} X, \quad y
\]
\note{petit schéma au bas de la page : $F' \ll F$ sur la première ligne, $F'_X \ll X$ au-dessous, reliés par des $\leq$ couchés ($F'_X \leq F'$, $X \leq F$) ; sous $X$, un $\gg$ surmonté d'un petit rond vers $x$, et $y$ à côté de $x$}

\page{82}
On veut remplacer la situation initiale par
\[
  F'_X \ll X, \qquad x \mathrel{\overset{\circ}{\ll}} X, \qquad y \qquad (\text{i.e. on } F \in \mathcal{M},\ G \in \mathcal{L})
\]
\note{sous $X$, le signe vertical vers $x$ est surchargé et raturé}
en utilisant \struck{deux} plus l'existence d'un $x \in \mathrm{omb}(X)^{\circ}$ tel que $x \mathrel{\overline{\parallel}} y$. Supposons qu'on puisse trouver $z \in \mathrm{omb}(X)$ telles que $z \parallel F'_X$. On aura $z \ll X \leq F$ donc $z \ll F$, prouvons que $z \parallel$ \struck{\ill{}} $F'$ (dans la figure, on prendra $G = z$).\note{la lettre $z$ est écrite chaque fois d'un trait barré, comme un $\mathcal{Z}$ ; « $G = z$ » : un accent ou prime après le $z$ n'est pas certain}

\struck{Mais en fait, on n'a pas \ill{} At 9 \uncertain{il suffit de} vérifier que $\forall\, z' \in \mathcal{L}$, $z' \ll F'$, on ait $z' \parallel z$, si $z' \ll F'_X$, cela résulte de $F'_X \parallel z$.} \uncertain{Cela ne résulte pas} \ill{}\note{passage barré d'un long trait et de traits obliques ; un cadre à gauche entoure « Mais en fait », « $z' \parallel z$ » et un mot commençant par « Di », lus en partie}

\emph{\uncertain{Le cas}} suivant \uncertain{suffit}
\[
  F' \ll F, \qquad F'_G \leq F', \qquad G \leq F, \qquad F'_G \ll G, \qquad H \ll G
\]
\drawing{schéma : $F' \ll F$ au-dessus de $F'_G \ll G$, reliés par des $\leq$ couchés ; sous $G$, un $\gg$ couché vers $H$}

Pour que \struck{$H$ soit disjoint de} l'on ait $H \parallel F'$, il (faut et) suffit que $H \parallel F'_G$.

(On l'appliquera : $G = X$, $H = z$) Supposons en effet $H \parallel F'_G$, prouvons $H \parallel F'$, donc que
\[
  x' \in \mathrm{omb}\, F',\ x \in \mathrm{omb}\, H \Longrightarrow x \parallel x' .
\]
\marginal{NB OPS $H = x$, $F' = x'$}\note{note encadrée à droite ; « $H = x$ » est écrit sur un premier signe biffé}
Or on sait par hyp.\ si $x \in \mathrm{omb}\, F'_G = \mathrm{omb}\, F' \cap \mathrm{omb}\, G$. Si $x \in \mathrm{omb}\, F' \smallsetminus \mathrm{omb}\, F'_G$

\page{83}
Considérons
\[
  x' \ll F' \ll F, \qquad x'_G \ll F'_G \ll G, \qquad x \ll G
\]
\drawing{schéma : deux lignes $x' \ll F' \ll F$ et $x'_G \ll F'_G \ll G$, reliées par trois $\leq$ couchés ($x'_G \leq x'$, $F'_G \leq F'$, $G \leq F$) ; sous $G$, un $\gg$ couché vers $x$}

Comme $x'_G \leq x'$, et $x' \in \mathcal{L}$, on aura $x'_G = x'$ ou $x'_G = \emptyset$, et comme $x' \not\ll F'_G$ par hyp., on a $x'_G = \emptyset$. Mais cela signifie aussi que \struck{$\mathrm{omb}(F$} $x' \in \mathrm{omb}(F) \smallsetminus \mathrm{omb}(G)$, donc $x$ et $x'$ appartiennent à deux strates différentes de $\mathrm{omb}(F)$, donc $x \parallel x'$ \struck{\ill{}} \ill{} At 10.\note{« $x' \not\ll F'_G$ » : le $\ll$ est barré d'un trait ; entre « $x \parallel x'$ » et « At 10 », quelques mots écrits en biais et biffés}

Nous allons donc poser

\marginal{NB At 8, At 9, At 13 sont des axiomes de type « beaucoup de lieux »}\note{note oblique dans la marge droite, en face de ce qui suit}

\emph{At 13} \struck{Soient} Soient $F \ll X$, avec $X \in \mathcal{M}$. \struck{Si} Supposons qu'il existe $G \in \mathfrak{F}$, $G \neq \emptyset$, $G \parallel F$, $G \mathrel{\overline{\parallel}} X$ (NB on \struck{pourrait} \ill{} pourra donc prendre $G \in \mathcal{L}$). Alors il existe $x \in \mathrm{omb}\, X$, tel que \struck{$F \parallel$} $x \parallel F$ (\ill{} $x \mathrel{\overline{\parallel}} X$).

(Cela bien \ill{} \ill{} dans

\marginal{« Axiome de \ill{} » (dans le critère d'existence \ill{}) ; Avec la notation \ill{} de \ill{} abstrait $(\mathfrak{F}, \leq)$ ou $(\mathcal{M}, \leq)$, l'hypothèse signifie aussi \ill{} ; At 13 dit que \ill{} $F \neq$ \ill{} il existe \ill{}}\note{longue note oblique dans la marge gauche, de At 13 au corollaire, lue par fragments}

\emph{Corollaire} Soit $F' \ll F$. Conditions équivalentes :
\begin{itemize}
  \item[(i)] $\forall\, G \in \mathfrak{F}$, $G \parallel F' \Longrightarrow G \parallel F$ \add{(donc $G \parallel F' \Leftrightarrow G \parallel F$)}
  \item[(ii)] \struck{$\forall\, x$} Pour tout $x \in \mathrm{omb}(X)$, on a $x \mathrel{\overline{\parallel}} F'$
\end{itemize}
\note{« $\mathrm{omb}(X)$ » est sur la page ; la démonstration de la page 84 se sert de $x \ll F$, i.e.\ $x \in \mathrm{omb}(F)$}

\page{84}
\begin{itemize}
  \item[(iii)] \struck{Il n'existe pas} \add{Pour tout $\overline{F'}$} \struck{\ill{}}, avec
  \[
    F' \leq \overline{F'} \ll F
  \]
  on a $F' = \overline{F'}$, i.e.\ $F'$ est maximal parmi les raffinements de $F$, ordonnés par $\leq$
  \item[(iv)] \struck{Pour tout $F'$ avec $F' \ll \overline{F'} \ll F$ on a $F' = \overline{F'}$, i.e.\ $F'$ est maximal parmi les raffinements de $F$, ordonnés par $\ll$.}
\end{itemize}
\note{(iv) est barré de grands traits en zigzag}

\emph{Dém.} (i) $\Leftrightarrow$ (ii) n'est autre que la prop.\ (conditionnée par At 13), compte tenu de la remarque qui la suit.

\struck{(iii) $\Rightarrow$ (ii) est évident} i.e.\ $\overline{\text{(ii)}} \Rightarrow \overline{\text{(iii)}}$. En effet (iii) $\Longrightarrow$ (ii) \struck{\ill{} (ii) n'était \ill{}} si \struck{\ill{}} on a $x \in \mathrm{omb}\, F$, avec $x \parallel F$,\note{« $\overline{\text{(ii)}}$ », « $\overline{\text{(iii)}}$ » : les numéros surlignés, pour les négations} considérons
\[
  F' \vee x = \overline{F'}
\]
prouvons que l'on a $\overline{F'} \ll F$ (ce qui montrera $\overline{\text{(iii)}}$, car $F' \nleq \overline{F'}$). Or \struck{il suffit donc de prouver cela} par At 4 cela se prouve strates par strates, \uncertain{on y gagne}.\note{« $x \parallel F$ » : on attend $x \parallel F'$. Le signe de « $F' \nleq \overline{F'}$ » est un $\leq$ barré ; on attend $F' \neq \overline{F'}$}

\struck{\emph{Prop.} Soit $F$ figure,} Il reste à prouver que (ii) $\Longrightarrow$ (iii). \struck{Prouvons d'abord (ii) $\Rightarrow$ (iii).} Si dans (iii), on avait $F' \neq \overline{F'}$, soit $X \in \widetilde{\overline{F'}} \smallsetminus \widetilde{F'}$, $x \in \mathrm{omb}(X)^{\circ}$, alors $x \ll \overline{F'} \ll F$ donc $x \ll F$, et $x \parallel F'$, contrairement à (ii). \struck{Ainsi (i) (ii) (iii) sont équivalentes et impliquées par (iv).} Mais impliquent-elles (iv) ?

\page{85}
\emph{Définition} Quand ces conditions sont satisfaites, on dit que $F'$ est une \emph{subdivision} de $F$, ou $F$ une \emph{surdivision} (ou une \emph{\uncertain{agrégation}}) de $F'$, et on écrit
\[
  F' \preccurlyeq F .
\]
\note{le signe est un $\leq$ dont le trait inférieur est ondulé, comme au lot 1 (page 7), où $\preccurlyeq$ note déjà la subdivision. Dans la première mouture (GF IV, dossier 156-4, lot 5, p. 84), $\preccurlyeq$ est une relation primitive dont se déduit $\ll$ ; ici la subdivision est définie à partir de $\ll$ et de la disjonction}

\marginal{NB \ill{} At 13, on \ill{} définition \ill{} $F' \ll F$ \ill{} ; NB \ill{} $F' \preccurlyeq F$ \ill{} surjectif \ill{} ; Plus généralement si $F'' \ll F' \ll F$ et \ill{} ; \emph{Mieux} : $F'' \ll F' \ll F$, ($F'' \preccurlyeq F$) $\Rightarrow$ $F' \preccurlyeq F$ \ill{}}\note{longues notes obliques dans la marge gauche, sur toute la moitié supérieure, lues par fragments}

\emph{Proposition} La relation $F' \preccurlyeq F$ dans $\mathfrak{F}$ est une relation d'ordre.

Elle est évidemment réflexive. Prouvons
\[
  F'' \preccurlyeq F' \preccurlyeq F \Longrightarrow F'' \preccurlyeq F .
\]
Soit \struck{donc} $G \in \mathfrak{F}$, supposons $G \parallel F''$, prouvons $G \parallel F$. En effet, on a $G \parallel F'$ car $F'' \preccurlyeq F'$, donc $G \parallel F$ car $F' \preccurlyeq F$, OK.

\emph{Proposition} Considérons
\[
  F' \ll F, \qquad F'_G \leq F', \qquad G \leq F, \qquad F'_G \ll G .
\]
\note{même schéma carré que page 82 ; près de « Proposition », de petits signes ($F''$, $\succcurlyeq$) mêlés à la note de marge}
Si $F' \preccurlyeq F$, alors $F'_G \preccurlyeq G$.

En effet, \struck{soit $K \in \mathfrak{F}$, avec $x$} il suffit de voir que pour \struck{$\forall\, x \in \mathrm{omb}$} tout $x \in \mathrm{omb}(G)$, on a $x \mathrel{\overline{\parallel}} F'_G$. Mais si on avait $x \parallel F'_G$, on aurait $x \parallel F'$ par le lemme p.~82, contrairement à $F' \preccurlyeq F$, cqfd.

\emph{Corollaire} Soit $F' \ll F$, et $F_i \leq F$ ($i \in I$) tels que $\operatorname{Sup} F_i = F$. Soit $F'_i = F' | F_i$ le raffinement de $F_i$ induit. Alors $F' \preccurlyeq F$ ssi $F'_i \preccurlyeq F_i$ $\forall\, i \in I$.

\page{86}
Nécessité résulte de la prop. Pour la suffisance, soit $x \in \mathrm{omb}\, F$, prouvons qu'on ne peut avoir $x \parallel F'$. Comme $\mathrm{omb}\, F = \bigcup_i \mathrm{omb}(F_i)$, $\exists\, i \in I$ t.q.\ $x \in \mathrm{omb}\, F_i$. Mais alors, comme $F'_i \preccurlyeq F_i$ on a $x \mathrel{\overline{\parallel}} F'_i$, ce qui par le lemme signifie $x \parallel F'$, cqfd.\note{« signifie $x \parallel F'$ » : la barre, qu'on attend, n'est pas sur la page}

\emph{Corollaire 2} Recollement des subdivisions.
\note{c'est, dans la première mouture, l'axiome C8 « Axiome de recollement des subdivisions » (GF IV, p. 86 ; dossier 156-4, lot 5)}

\emph{At 14} Soient $G \leq F$, $G' \preccurlyeq G$. Alors $\exists$ subdivision $F'$ de $F$ qui induise $G'$.
\note{sous $G$, un $\succcurlyeq$ couché vers $G'$. At 14 est marqué d'un trait vertical dans la marge}

\marginal{\uncertain{Je voudrais} \ill{} ?}\note{note oblique dans la marge gauche, en face de At 14}

Par le recollement des subdivisions, on doit être ramené (par une construction transfinie \ill{}, si $\widetilde{F}$ infini) au cas où
\[
  F = X \in \mathcal{M} .
\]

\emph{Remarque} On aimerait avoir un choix canonique, ou une subdivision $F'$ \uncertain{la moins fine} \struck{\ill{}} (maximale) parmi celles qui ont la propriété voulue, comme dans le cas ensembliste. Ce serait une variante, qu'on pourrait énoncer ainsi :

\page{87}
\emph{At \struck{\ill{}} (fort) 14 bis} Sous les conditions précédentes, \struck{considérons} Considérons \struck{Multomb}
\[
  \widetilde{G} \subset \widetilde{F}, \qquad \widetilde{G'} \longrightarrow \widetilde{G}
\]
\note{la flèche monte de $\widetilde{G'}$ vers $\widetilde{G}$, sous l'inclusion $\widetilde{G} \subset \widetilde{F}$ ; « (fort) » est écrit par-dessus un mot biffé}
alors on peut trouver $F' \preccurlyeq F$ induisant $G'$, de telle façon que dans le diagramme d'ens.\ ordonnés

\begin{tikzcd}
  \widetilde{G'} \arrow[r, hook] \arrow[d, "\varphi|\widetilde{G'}"'] & \widetilde{F'} \arrow[d, "\varphi"] \\
  \widetilde{G} \arrow[r, hook] & \widetilde{F}
\end{tikzcd}

\note{carré marqué « cart » au centre ; l'étiquette de la flèche de gauche, lue $\varphi|\widetilde{G'}$, est douteuse}
on ait
\[
  \widetilde{F'} \smallsetminus \widetilde{G'_{+}} \xrightarrow{\ \sim\ } \widetilde{F} \smallsetminus \widetilde{G},
\]
\note{au-dessus de « $\widetilde{F'} \smallsetminus \widetilde{G'_{+}}$ » : « sous-ens.\ fermé » ; au-dessus de la flèche : « iso d'ens.\ ordonnés ». L'indice lu $+$ est douteux}
et pour $\xi \in \widetilde{G'_{+}}$, $\eta \in \widetilde{F'}$, on ait $\xi \leq \eta$ ssi $\varphi^{*}(\xi) \leq \varphi^{*}(\eta)$ [ce qui détermine l'ens.\ ordonné $\widetilde{F'}$ et les diagrammes précédents, à iso unique près].\note{l'exposant de $\varphi$, une petite étoile ou un dièse, est rendu $\ast$}

\marginal{\emph{22. Juin}}\note{date de sa main, soulignée, dans la marge gauche, en face du paragraphe qui suit}

J'ai l'impression qu'on doit pouvoir prouver qu'un tel $F'$ est unique, comme la plus grande des subdivisions de $F$ induisant $G'$ sur $G$. Dans le cas où $\mathfrak{F}$ est à figures finies (pour éviter des inductions transfinies \ldots), \struck{cette} cette condition est équivalente à :

\marginal{unicité \ill{} ?}\note{note oblique dans la marge gauche}

\emph{At (fort) 14 ter} : Soit $X \in \mathcal{M}$, et $F' \preccurlyeq \partial X$. Alors il existe un unique $X' \in \mathcal{M}$ tel que $X' \preccurlyeq X$, $\partial X' = F'$.
\note{l'énoncé est marqué d'un double trait vertical dans la marge}

\page{88}
Nous imposerons du moins cette condition pour les ateliers \emph{stricts} (cf plus bas). Mais on notera que cette condition n'est pas vérifiée dans le cas d'ateliers \struck{\ill{}} simpliciaux \add{stricts}, par exemple [triangle] est un contre-exemple.\note{un petit triangle dessiné, un point marqué près d'un sommet, précédé d'un « ! »}

(Atelier simplicial : \struck{tout $X \in \mathcal{M}$} $\forall\, X \in \mathcal{M}$ est un « simplexe » \ill{} \struck{toute} \ill{} partie \add{non vide} \struck{\ill{}} de $\widetilde{X}$ \ill{} \ldots\ \emph{\ill{}} \ill{} $\widetilde{X}$ \struck{soit} finis, \struck{\ill{} et tout $X$ soit fini.} \ill{} de $\mathcal{M}$ \ill{} \ldots\ par exemple : $\bigcirc$ )\note{parenthèse très corrigée : plusieurs lignes serrées, des insertions et un passage barré de plusieurs traits ; lue en partie seulement. Elle s'achève sur un petit cercle dessiné}

\marginal{Atelier simplicial strict : simplicial et $X, Y$ compatibles $\Rightarrow$ $X \cap Y \in \mathcal{M}$, i.e.\ $(X, Y)$ \ill{} dans $\mathcal{M}$ \ill{}}\note{note oblique dans la marge gauche, lue en partie}

Pour le moment, je voudrais
\begin{itemize}
  \item[1°)] Expliciter les axiomes dans le cas d'un atelier ensembliste
  \item[2°)] Voir leur stabilité \ill{} \ill{} stabilité par passage à un sous-atelier
  \item[3°)] Étudier l'opération de \emph{saturation} d'un atelier (\ill{} \ill{} au cas d'un atelier « modéré »).
\end{itemize}

Nous avons vu (p.~\uncertain{55}-57) qu'un atelier ensembliste est défini par un ens.\ $\mathcal{L}$ et une partie \struck{\ill{}}
\[
  \mathfrak{F} \subset \mathrm{Fig}(\mathcal{L})
\]
satisfaisant les conditions Atens 1, 2, 3. (On a \uncertain{déjà vu} qu'on a automatiquement At 1 à At 8. Il faut examiner les autres axiomes. Pour les exemples et contre-exemples on pensera surtout au cas $\mathcal{L}$ fini, et \ill{} on peut se représenter au \uncertain{partie} \ldots\note{« Atens » est suivi d'un petit signe en exposant ; la ligne finale est écrite serrée au bas du feuillet}

\page{89}
\struck{At 9 J'ai pensé \ill{} au cas}

\ill{} d'une donnée de $\mathcal{M}$ ens.\ de figures \ill{} dans $\mathcal{L}$, avec relation $\underset{\mathcal{M}}{\between}$, satisfaisant Atens* 1 -- Atens* 3. Si on prend p.ex.
\[
  \mathcal{M} = \lbrace \lbrace \lbrace x \rbrace \rbrace \mid x \in \mathcal{L} \rbrace ,
\]
\note{le premier départ de l'accolade est surchargé, et un mot est biffé sous « $x \in \mathcal{L}$ ». Le sigle lu « Atens* » (au lot 4, page 61 : « Atens* 4 ») se présente ici comme « At ens » suivi d'un signe proche de « ff »}
donc la donnée se réduit strictement à celle d'une relation \add{symétrique réflexive} $x \between y$ dans $\mathcal{L}$ (les conditions Atens* 1--3 étant trivialement satisfaites).

At 9 dans ce cas-là \struck{signifie que} \add{est \uncertain{automatiquement} vérifié}. Dans \struck{le} cas général, il se traduit par

\textbf{\struck{Atens* 4}} (*) \textbf{Atens* 4} Si $X, Y \in \mathcal{M}$ et \add{$|X| \cap |Y| = \emptyset$. Alors $X \underset{\mathcal{M}}{\between} Y$ ssi} $\forall\, x \in |X|$, $y \in |Y|$, $x \underset{\mathcal{M}}{\between} y$. \struck{\ill{}}
\note{la fin de l'énoncé est ajoutée au-dessus de la ligne et rattachée par un trait ; le sigle est surchargé et biffé, puis récrit}

Mais compte tenu de \emph{Atens} 2, cela équivaut à :

\emph{Atens 2} (ou Atens* 4) Si $X, Y \in \mathcal{M}$ et $|X| \cap |Y| = \emptyset$, alors $X \underset{\mathcal{M}}{\between} Y$ (donc $X \parallel Y$).

\marginal{En fait, on a donc $X \parallel Y$ ssi $|X| \cap |Y| = \emptyset$}\note{note oblique dans la marge gauche}

Atens $9_{\Sigma}$ se traduit par

\emph{Atens $2_{\Sigma}$} (ou Atens* $4_{\Sigma}$) \struck{Soient} Soient $X, Y \in \mathcal{M}$, \add{(\ill{} fig.\ ens.\ \ill{})} soit $K$ l'ens.\ des strates communes à $X$ et $Y$. Pour que $X \underset{\mathcal{M}}{\between} Y$, il faut et il suffit que $|X| \cap |Y| = K$ (i.e.\ $X \underset{\mathrm{ens}}{\between} Y$) \emph{et} $K \in \Sigma$.\note{l'énoncé est marqué d'un trait vertical dans la marge}

\page{90}
Le cas où $\Sigma =$ tous les ens.\ ordonnés (axiome At 9') ceci équivaut à dire que $\mathfrak{F}$ est un sous-atelier \emph{spécieux} de \struck{$\mathfrak{F}$} $\mathrm{Fig}(\mathcal{L})$. Mais \struck{il} on se rappellera que ce ne sera pratiquement jamais le cas pour les ateliers \struck{\ill{}} simpliciaux.

At 10 automatique par Atens (2).

\emph{At \struck{\ill{}}} Axiome de stabilité par Sup.\ filtrantes $\operatorname{Sup} F_i$, quand les $F_i \ll F$ et la famille $F_i$ est « loc.\ finie près de $F$ ».\note{le numéro de l'axiome est biffé et n'est pas remplacé ; c'est l'axiome des Sup filtrants de la page 80 (At 11 bis). La page porte « $\operatorname{Sup} F'_i$ », le prime douteux}

\marginal{\ill{} de \ill{} « très \ill{} »}\note{note oblique dans la marge gauche}

Cet axiome est vide dans le cas où les figures $\in \mathfrak{F}$ sont \struck{de t.f.\ finies} de t.f.\ (car alors la famille $(F_i)$ \ill{} $=$ \uncertain{stationnaire} \ldots) At 12 (ou l'axiome donné \ill{} At $9_{\Sigma}$ i.e.\ Atens $2_{\Sigma}$)\note{« At 12 » est écrit sous un sigle biffé ; la parenthèse est lue en partie}

\emph{Atens 4} (ou Atens* 4) Soient \struck{$X, Y \in \mathcal{M}$}, $X' \ll X$, $Y' \ll Y$ dans $\mathcal{M}$, \struck{$K = X \cap Y$} avec $X \underset{\mathcal{M}}{\between} Y$, $K = X \cap Y$. Si $X'_K \underset{\mathcal{M}}{\between} Y'_K$ (i.e.\ pour tout \struck{\ill{}} $X'' \in \widetilde{X'}$ telle que $X'' \ll K$ (i.e.\ $X'' \ll Y$) et tout $Y'' \in \widetilde{Y'}$ telle que $Y'' \ll K$ (i.e.\ $Y'' \ll X$) on a $X'' \underset{\mathcal{M}}{\between} Y''$), alors $X \underset{\mathcal{M}}{\between} Y$.\note{« alors $X \between Y$ » est sur la page ; c'est la traduction de At 12 (lot 4, page 76), dont la conclusion est $X' \between Y'$. Le chiffre de « Atens* 4 » est surchargé}

\marginal{Résulte de Atens $2_{\Sigma}$}\note{note oblique dans la marge gauche, en face de Atens 4}

Grâce à Atens 2, on voit que \struck{At} la condition sur $F \ll X$ dans At 13 signifie que $|F| \subsetneq |X|$, et \struck{\ill{}} que At 13 est toujours vérifié.

\page{91}
La condition At 14 (peut-être \struck{\ill{}} condition) semble \struck{plus} un peu \uncertain{discordante}, traduite en termes de sous-ateliers d'un $\mathrm{Fig}(\mathcal{L})$ --- ce n'est plus dans le \ill{} d'une condition de stabilité. \ill{}, l'axiome At 14 bis (\uncertain{quantitatif}) est plus sympathique :

\emph{Atens fort 5} : Soit $X \in \mathcal{M}$, et $F' \preccurlyeq \partial X$. Considérons l'unique $F' \in \mathcal{M}$. figure $F \in \mathrm{Fig}(\mathcal{L})$ qui soit élémentaire \struck{---} et telle que $\partial F = F'$. Alors $F \in \mathcal{M}$.\note{« $F' \in \mathcal{M}$. » est seul sur sa ligne, sans suite ; la phrase reprend au-dessous. L'énoncé est marqué d'un trait vertical dans la marge}

Mais ce sera pas vrai pour les ateliers simpliciaux !

\note{un court trait horizontal sépare ce qui suit}

Passons à la situation d'un atelier $\mathfrak{F}$, et d'un sous-atelier $\mathfrak{F}' \subset \mathfrak{F}$ (cf conditions Sous at 1, 2 p.~\uncertain{47})\note{la page 47 est au lot 3}

Les conditions At 1 à At 7 sur \struck{\ill{}} $\mathfrak{F}$ se transportent automatiquement à $\mathfrak{F}'$ ! Pour les conditions impliquant $\mathcal{L}$ (At 8, 9, $9_{\Sigma}$, 10) il faut distinguer soigneusement $\mathcal{L}$ et $\mathcal{L}'$. Évidemment
\[
  \mathcal{L} \cap \mathfrak{F}' \subset \mathcal{L}' .
\]

\page{92}
On voudrait qu'il y ait égalité, i.e.\ que tout élément minimal de $\mathfrak{F}'$ \add{(pour $\ll$)} \struck{\ill{}} (\ill{}) soit minimal dans $\mathfrak{F}$. On posera

\emph{Définition} On dit que le sous-atelier $\mathfrak{F}'$ de $\mathfrak{F}$ est \emph{épais} (pour les lieux) si $F \in \mathfrak{F}' \Rightarrow \mathrm{omb}(F) \subset \mathfrak{F}'$ i.e.\
\[
  F \in \mathfrak{F}',\ x \ll F \ (x \in \mathcal{L}) \Longrightarrow x \in \mathfrak{F}' .
\]
Il est dit \struck{\ill{}} \add{\uncertain{quasi-dense} si} $\forall\, X \in \mathcal{M}' = \mathfrak{F}' \cap \mathcal{M}$, $\exists\, x \in \mathcal{L} \cap \mathfrak{F}'$ avec $x \mathrel{\overset{\circ}{\ll}} X$ i.e.\ $x \in \mathrm{omb}(X)^{\circ}$.

Cela équivaut à :

\emph{Sous at} (épais) 3 \quad $\forall\, X \in \mathcal{M}'$, $x \in \mathcal{L}$, t.q.\ $x \ll X$, $x \in \mathfrak{F}'$.

\struck{On dit} \add{ou} \struck{Corollaire}

\emph{Sous-at} (quasi-dense) 3' \quad $\forall\, X \in \mathcal{M}' = \mathfrak{F}' \cap \mathcal{M}$, $\exists\, x \in \mathcal{L} \cap \mathfrak{F}'$ avec $x \mathrel{\overset{\circ}{\ll}} X$.

Bien sûr, Sous at 3 $\Longrightarrow$ Sous at 3', et ces deux conditions impliquent At 8 pour $\mathfrak{F}'$, et la dernière implique déjà \struck{\ill{}}
\[
  \mathcal{L}' = \mathcal{L} \cap \mathfrak{F}' .
\]
\note{avant la formule, trois mots soulignés puis biffés}

Ceci dit, At 9 n'est pas nécessairement stable par passage à un « sous-atelier », même \uncertain{plein}. Pour être \uncertain{correct}, on voudra \add{\ill{}}

\emph{Sous-at 4} \quad Si $X, Y \in \mathcal{M}'$ et $X \parallel Y$ dans $\mathfrak{F}$, alors $X \parallel Y$ dans $\mathfrak{F}'$ ($\Leftrightarrow X \underset{\mathfrak{F}'}{\between} Y$)

\marginal{dans la \ill{} « propriété »}\note{note oblique dans la marge gauche, en face de Sous-at 4}

\page{93}
qui \ill{} implique, \struck{\ill{}} avec Sous-at 3 (atelier plein) que si $\mathfrak{F}$ satisfait At 9, $\mathfrak{F}'$ aussi (bien que dans Sous-at 4 il soit question des lieux).

Moyennant Sous-at 4, \add{et Sous-at 3}, \struck{\ill{}} At 10 pour $\mathfrak{F}'$ implique \struck{que} itou pour $\mathfrak{F}$.\note{ainsi sur la page ; le sens attendu, d'après le reste de la page, est de $\mathfrak{F}$ vers $\mathfrak{F}'$}

At 11 bis pour $\mathfrak{F}'$ est une propriété de stabilité \add{de $\mathfrak{F}'$ dans $\mathfrak{F}$} \ill{} automatique.\note{sous la ligne de At 11 bis est ajouté « (\ill{} At 11) »}

Idem pour At 12. Pour At 13, \struck{\ill{}} moyennant Sous-at 3, Sous-at 4, la validité dans $\mathfrak{F}$ l'implique dans $\mathfrak{F}'$. At 14 est \uncertain{discordante}, At 14 bis dans le \ill{} d'une condition de stabilité.

\note{un court trait horizontal clôt la page, le chapitre et le dossier}

\end{document}
